\documentclass[10pt,a4paper]{amsart}
\usepackage{amsmath,amssymb,amsthm,graphicx}
\usepackage{verbatim}
\newenvironment{reference}{\comment}{\endcomment}
%\newenvironment{reference}{}{}
\newenvironment{slogan}{\comment}{\endcomment}
\newenvironment{history}{\comment}{\endcomment}

% For commutative diagrams we use Xy-pic
\usepackage[all]{xy}

% We use 2cell for 2-commutative diagrams.
\xyoption{2cell}
% Theorem environments.
%
\theoremstyle{plain}
\newtheorem{theorem}[subsection]{Theorem}
\newtheorem{proposition}[subsection]{Proposition}
\newtheorem{lemma}[subsection]{Lemma}

\theoremstyle{definition}
\newtheorem{definition}[subsection]{Definition}
\newtheorem{example}[subsection]{Example}
\newtheorem{exercise}[subsection]{Exercise}
\newtheorem{situation}[subsection]{Situation}

\theoremstyle{remark}
\newtheorem{remark}[subsection]{Remark}
\newtheorem{remarks}[subsection]{Remarks}

\numberwithin{equation}{subsection}

% Macros
%
\def\lim{\mathop{\mathrm{lim}}\nolimits}
\def\colim{\mathop{\mathrm{colim}}\nolimits}
\def\Spec{\mathop{\mathrm{Spec}}}
\def\Hom{\mathop{\mathrm{Hom}}\nolimits}
\def\Ext{\mathop{\mathrm{Ext}}\nolimits}
\def\SheafHom{\mathop{\mathcal{H}\!\mathit{om}}\nolimits}
\def\SheafExt{\mathop{\mathcal{E}\!\mathit{xt}}\nolimits}
\def\Sch{\mathit{Sch}}
\def\Mor{\mathop{\mathrm{Mor}}\nolimits}
\def\Ob{\mathop{\mathrm{Ob}}\nolimits}
\def\Sh{\mathop{\mathit{Sh}}\nolimits}
\def\NL{\mathop{N\!L}\nolimits}
\def\CH{\mathop{\mathrm{CH}}\nolimits}
\def\proetale{{pro\text{-}\acute{e}tale}}
\def\etale{{\acute{e}tale}}
\def\QCoh{\mathit{QCoh}}
\def\Ker{\mathop{\mathrm{Ker}}}
\def\Im{\mathop{\mathrm{Im}}}
\def\Coker{\mathop{\mathrm{Coker}}}
\def\Coim{\mathop{\mathrm{Coim}}}

% Boxtimes
%
\DeclareMathSymbol{\boxtimes}{\mathbin}{AMSa}{"02}

%
% Macros for moduli stacks/spaces
%
\def\QCohstack{\mathcal{QC}\!\mathit{oh}}
\def\Cohstack{\mathcal{C}\!\mathit{oh}}
\def\Spacesstack{\mathcal{S}\!\mathit{paces}}
\def\Quotfunctor{\mathrm{Quot}}
\def\Hilbfunctor{\mathrm{Hilb}}
\def\Curvesstack{\mathcal{C}\!\mathit{urves}}
\def\Polarizedstack{\mathcal{P}\!\mathit{olarized}}
\def\Complexesstack{\mathcal{C}\!\mathit{omplexes}}
% \Pic is the operator that assigns to X its picard group, usage \Pic(X)
% \Picardstack_{X/B} denotes the Picard stack of X over B
% \Picardfunctor_{X/B} denotes the Picard functor of X over B
\def\Pic{\mathop{\mathrm{Pic}}\nolimits}
\def\Picardstack{\mathcal{P}\!\mathit{ic}}
\def\Picardfunctor{\mathrm{Pic}}
\def\Deformationcategory{\mathcal{D}\!\mathit{ef}}
\usepackage{cleveref}
\crefname{originallemma}{original lemma}{original lemmas}
\Crefname{originallemma}{Original lemma}{Original lemmas}
\makeatletter
\newlabel{algebra-lemma-etale-quasi-finite}{{00U5}{source}{}{sourcepointer-1}{}}
\newlabel{algebra-lemma-etale-quasi-finite@cref}{{[originallemma][0][]00U5}{[1][0][]source}{}{}{}}
\makeatother
\crefname{originallemma}{original lemma}{original lemmas}
\Crefname{originallemma}{Original lemma}{Original lemmas}
\makeatletter
\newlabel{algebra-lemma-map-between-etale}{{00U7}{source}{}{sourcepointer-2}{}}
\newlabel{algebra-lemma-map-between-etale@cref}{{[originallemma][0][]00U7}{[1][0][]source}{}{}{}}
\makeatother
\crefname{originallemma}{original lemma}{original lemmas}
\Crefname{originallemma}{Original lemma}{Original lemmas}
\makeatletter
\newlabel{algebra-lemma-silly}{{00DS}{source}{}{sourcepointer-3}{}}
\newlabel{algebra-lemma-silly@cref}{{[originallemma][0][]00DS}{[1][0][]source}{}{}{}}
\makeatother
\crefname{originallemma}{original lemma}{original lemmas}
\Crefname{originallemma}{Original lemma}{Original lemmas}
\makeatletter
\newlabel{algebra-lemma-strict-henselization-different}{{04GW}{source}{}{sourcepointer-4}{}}
\newlabel{algebra-lemma-strict-henselization-different@cref}{{[originallemma][0][]04GW}{[1][0][]source}{}{}{}}
\makeatother
\crefname{originalproposition}{original proposition}{original propositions}
\Crefname{originalproposition}{Original proposition}{Original propositions}
\makeatletter
\newlabel{algebra-proposition-fppf-open}{{00I1}{source}{}{sourcepointer-5}{}}
\newlabel{algebra-proposition-fppf-open@cref}{{[originalproposition][0][]00I1}{[1][0][]source}{}{}{}}
\makeatother
\title{A ring over which every faithfully flat etale algebra has a retraction has strictly henselian local rings at maximal ideals}
\author{The Stacks Project Authors}
\date{stacks-project:proetale.tex:local-snapshot-20261012; GNU FDL 1.2 or later}
\begin{document}
\maketitle
\noindent\textit{Editorial note.} All delivered mathematical text is the author's own, copied from the pinned Stacks Project chapter file proetale.tex; no statement or proof step was rewritten, reordered or removed.\par\medskip
\noindent\textit{Editorial note.} The counted proof is the author's complete argument: the reduction to an isomorphism of the local rings through the description of the strict henselization, the separation of the primes over the maximal ideal, the construction of a faithfully flat etale map with a retraction via the openness of the image, and the closing flatness argument.\par\medskip
\noindent\textit{Editorial note.} The label lemma-have-sections-strictly-henselian names the counted statement, and the proof environment immediately following it in the source is delivered verbatim as its proof.\par\medskip
\noindent\textit{Editorial note.} The counted statement is carried by the author's own lemma environment in the chapter source.\par\medskip
\noindent\textit{Editorial note.} Pre-open gate: python operations/agent-authoring/tools/precard\_gate.py --source proetale.tex --label lemma-have-sections-strictly-henselian --start 1660 --end 1700 --local-id 3580 returned RESULT=OK (L3580 free, occupancy FREE, public interval clear). Candidate taken from sources/stacks-proetale/candidates-halfA.json.\par\medskip
\noindent\textit{Editorial note.} The statement and proof are copied from the pinned chapter file; the cited inputs are kept as external pointers with their canonical Stacks tags.\par\medskip
\noindent\textit{Allowed standard prerequisites (editorial).}\begin{itemize}
\item Description of the strict henselization of a local ring. Algebra, the lemma describing the strict henselization of a local ring as the limit over the etale neighbourhoods of a prime, used to reduce to an isomorphism of local rings. Supplies the reduction of the required statement to showing that the local map A\_m -> B\_q is an isomorphism.
\item Quasi-finiteness of etale ring maps and the finiteness of the primes over a maximal ideal. Algebra, the lemma that an etale ring map is quasi-finite, and the elementary lemma allowing one to separate finitely many primes by an element. Supply the fact that the finitely many primes over m are maximal and admit no specializations, and the existence of an element separating them.
\item Openness of the image of a flat finitely presented morphism. Algebra, the proposition that the image of a flat and finitely presented morphism of affine schemes is open. Supplies the openness of the image U of Spec(B) -> Spec(A), hence the element f with Spec(A) = U union D(f).
\item Flat local homomorphisms of local rings are injective, and maps between etale algebras are flat. Algebra, the lemma that a flat local homomorphism of local rings is injective, and the lemma that a map between etale algebras over a ring is flat. Supply the conclusion that the retraction A\_p -> B\_q -> A\_p forces the two local maps to be isomorphisms.
\end{itemize}\medskip
\ifdefined\hypertarget\hypertarget{sourcepointer-1}{}\fi
\noindent\textit{Source reference.} Original lemma 00U5: sources/stacks-chow/tags.txt, source lines 1099--1099. Prerequisite pointer invoked at proof line(s) 1678; it names a standard fact of a sibling chapter that the counted argument uses, and every step of the counted proof built on it is retained verbatim in the delivered body.\par\medskip
\ifdefined\hypertarget\hypertarget{sourcepointer-2}{}\fi
\noindent\textit{Source reference.} Original lemma 00U7: sources/stacks-chow/tags.txt, source lines 1101--1101. Prerequisite pointer invoked at proof line(s) 1692; it names a standard fact of a sibling chapter that the counted argument uses, and every step of the counted proof built on it is retained verbatim in the delivered body.\par\medskip
\ifdefined\hypertarget\hypertarget{sourcepointer-3}{}\fi
\noindent\textit{Source reference.} Original lemma 00DS: sources/stacks-chow/tags.txt, source lines 508--508. Prerequisite pointer invoked at proof line(s) 1681; it names a standard fact of a sibling chapter that the counted argument uses, and every step of the counted proof built on it is retained verbatim in the delivered body.\par\medskip
\ifdefined\hypertarget\hypertarget{sourcepointer-4}{}\fi
\noindent\textit{Source reference.} Original lemma 04GW: sources/stacks-chow/tags.txt, source lines 5847--5847. Prerequisite pointer invoked at proof line(s) 1672; it names a standard fact of a sibling chapter that the counted argument uses, and every step of the counted proof built on it is retained verbatim in the delivered body.\par\medskip
\ifdefined\hypertarget\hypertarget{sourcepointer-5}{}\fi
\noindent\textit{Source reference.} Original proposition 00I1: sources/stacks-chow/tags.txt, source lines 663--663. Prerequisite pointer invoked at proof line(s) 1685; it names a standard fact of a sibling chapter that the counted argument uses, and every step of the counted proof built on it is retained verbatim in the delivered body.\par\medskip
\begin{lemma}
\label{lemma-have-sections-strictly-henselian}
Let $A$ be a ring such that every faithfully flat \'etale ring map
$A \to B$ has a retraction. Then every local ring of $A$ at a maximal
ideal is strictly henselian.
\end{lemma}

\begin{proof}
Let $\mathfrak m$ be a maximal ideal of $A$. Let $A \to B$ be an
\'etale ring map and let $\mathfrak q \subset B$ be a prime
lying over $\mathfrak m$. By the description of the strict henselization
$A_\mathfrak m^{sh}$ in
Algebra, Lemma \ref{algebra-lemma-strict-henselization-different}
it suffices to show that $A_\mathfrak m = B_\mathfrak q$.
Note that there are finitely many primes
$\mathfrak q = \mathfrak q_1, \mathfrak q_2, \ldots, \mathfrak q_n$
lying over $\mathfrak m$ and there are no specializations
between them as an \'etale ring map is quasi-finite, see
Algebra, Lemma \ref{algebra-lemma-etale-quasi-finite}.
Thus $\mathfrak q_i$ is a maximal ideal and we can find
$g \in \mathfrak q_2 \cap \ldots \cap \mathfrak q_n$, $g \not \in \mathfrak q$
(Algebra, Lemma \ref{algebra-lemma-silly}).
After replacing $B$ by $B_g$ we see that $\mathfrak q$
is the only prime of $B$ lying over $\mathfrak m$.
The image $U \subset \Spec(A)$ of $\Spec(B) \to \Spec(A)$ is
open (Algebra, Proposition \ref{algebra-proposition-fppf-open}).
Thus the complement $\Spec(A) \setminus U$ is closed
and we can find $f \in A$, $f \not \in \mathfrak p$ such that
$\Spec(A) = U \cup D(f)$. The ring map $A \to B \times A_f$
is faithfully flat and \'etale, hence has a retraction
$\sigma : B \times A_f \to A$ by assumption on $A$.
Observe that $\sigma$ is \'etale, hence flat as a map between \'etale
$A$-algebras (Algebra, Lemma \ref{algebra-lemma-map-between-etale}).
Since $\mathfrak q$ is the only prime of $B \times A_f$ lying
over $A$ we find that $A_\mathfrak p \to B_\mathfrak q$ has
a retraction which is also flat. Thus
$A_\mathfrak p \to B_\mathfrak q \to A_\mathfrak p$
are flat local ring maps whose composition is the identity. Since
a flat local homomorphism of local rings is injective we conclude these
maps are isomorphisms as desired.
\end{proof}

\end{document}
